%% Sect. 6 -- from augur_AMT_§6 draft v1 (2026-09-20, later revisions to 2026-09-27)
%% The draft's "6.1a" is typeset as subsubsection 6.1.1 (label sec:6.1a).
\section{What should be reported, and what should be changed}
\label{sec:6}

Section~\ref{sec:4} established that the fit-reported uncertainty omits terms that are
several times larger than itself, and Sect.~\ref{sec:5} that convergence does not
imply identifiability. Neither finding requires a new retrieval algorithm to
act on. This section states what a DOAS product should carry instead of a
single uncertainty, and what in the measurement schedule is worth changing.
Of the four recommendations, two are implemented in Augur and evaluated against
the campaign: the knot gate (Sect.~\ref{sec:6.3}) and the coverage fields
(Sect.~\ref{sec:6.4}). The reporting convention of Sect.~\ref{sec:6.1} is implemented in
part -- the termination state, the distance to the nearest knot and the identity of
the reference file are exported, while the build time and time convention of each
reference are specified for implementation -- and the question of block length in
Sect.~\ref{sec:6.2} is left open by the measurements reported here. Section~\ref{sec:6.5} names the errors that none
of these measures reaches.

\subsection{Three numbers per retrieval, not one}
\label{sec:6.1}

Every retrieved amount is reported with
\begin{enumerate}
\item the \strong{fit uncertainty}, from the covariance of the projected problem
   (Sect.~\ref{sec:2.5}) -- what the residual supports, and nothing more;
\item the \strong{structural term}, evaluated for that scan's own calibration
   context rather than as a campaign constant; and
\item a \strong{context flag} recording the state of the references the scan depends
   on, in the specific form set out below.
\end{enumerate}

The third item is the one that has to be specified rather than gestured at,
because its value is entirely in which fields it carries. We report, per
record and per reference type (zero air, reflectivity), the fields of Table~\ref{tab:context_fields}.

\begin{table}[t]
\caption{Fields of the context flag, reported per record and per reference type (zero air, reflectivity), and the reason for each.}
\label{tab:context_fields}
\begin{tabular}{>{\raggedright\arraybackslash}p{0.38\textwidth}p{0.55\textwidth}}
\tophline
field & why \\
\middlehline
seconds to the nearest knot, with the side (before/after) & the side distinguishes interpolation from extrapolation \\
length of the bracketing interval & an eleven-hour gap and a fifteen-minute one are not the same interpolation \\
in-range flag & extrapolation is a different operation and should be nameable in a filter \\
number of adjacent calibration blocks rejected by the quality gates & the gates' cost, which Sect.~\ref{sec:6.4} argues must be reported wherever the gates are applied \\
identifier of the reference file, and when it was built & two products made from different calibration vintages are not comparable \\
the time convention the reference was built under & a correction applied upstream makes every earlier derived product stale \\
\bottomhline
\end{tabular}
\end{table}

The last field is there because of Sect.~\ref{sec:5.4}, and it is worth being exact about
what it would and would not have caught. The heated-channel clock ran on local time until 29 May
and on UTC thereafter, so for twelve days of the campaign the raw time stamps were nine hours off the
convention of the time axis; that is a property of the raw record, and no
provenance field detects it. It was detected the way such things are detected
-- against an external time reference -- and a correction was added to the
raw-data reader two months after the deployment ended.

What the field catches is the failure that followed. The reflectivity knots
had been computed before the correction existed, and nothing rebuilt them
afterwards. From that moment the retrieval was combining spectra on the
corrected time axis with a reference on the superseded one, and the twelve
days of product made from that combination are wrong by more than their whole
structural budget. This is not a clock problem; it is a staleness problem, and
it is general. Any correction applied upstream invalidates every derived
product built before it, and a pipeline in which derived products do not
record the convention they were built under cannot tell that it has happened.

The mismatch was not silent in the data. It surfaced as a nine-hour run of
out-of-range records at the very start of the deployment -- the first
zero-air block at 02:42, the first reflectivity knot at 11:40 -- and \fieldnum{42.3}\,\% of
that day extrapolated. The distance and in-range fields would have
carried that anomaly into the product on the first run after the correction,
rather than two months later. They would not have explained it: identifying
the cause needed the build time of the reflectivity file and the convention it
was built under. Detection and diagnosis are separate functions and need
separate fields.

There is also a cheaper protection than reporting, and it belongs in the same
subsection because it is the same information used earlier. A reference
carrying a time-convention tag can be checked against the tag on the spectra
at load time, and a mismatch refused rather than reported. Reporting turns a
silent error into a visible one; an assertion prevents it. We recommend the
assertion where the tags exist and the report in every case, because the
report is also what makes the error interpretable a year later, when the
person reading the product is not the person who ran it.

The second of these cannot be obtained cheaply, and it is worth saying why.
The amplification factor $I_0/(I_0-I)$, which governs how an error in the
zero-air reference enters the extinction, varies by a factor of forty between
records and is available without any fitting; it is a natural candidate for a
closed-form structural term. It does not work. Regressed against the measured
response it explains between \fieldnum{0.0}\,\% and \fieldnum{9}\,\% of the variance, with the wrong
sign, because the polynomial background absorbs the smooth part of the
perturbation and what reaches the gas columns is set by shape rather than
amplitude. Normalising the response by the retrieved amount does not help
either: the robust coefficient of variation is \fieldnum{0.80}--\fieldnum{1.09} in relative terms
against \fieldnum{0.83}--\fieldnum{1.11} in absolute. The structural term has to be obtained by
running the perturbation. That is available in post-processing and not in a
real-time monitor, and the two implementations differ accordingly.

The second and third exist because the first is conditional. A fit
uncertainty is a statement about the residual of a model assumed correct; it
is not a statement about the inputs to that model. Separating the three keeps
that distinction visible to whoever uses the product, instead of resolving it
silently into one number that is too small.

\subsubsection{Which remedy a term needs is readable from its shape}
\label{sec:6.1a}

The three structural terms call for three different actions, and the shape
statistic $s$ of Sect.~\ref{sec:4.3} says which. A term with $s \approx 1$ is Gaussian: its
error is spread evenly over all records and can only be reduced by a better
model. A term with $s \approx 0.7$--1 is dominated by measurement noise on the
references and responds to averaging. A term with $s \ll 0.7$ is concentrated
in a few records, where averaging is useless and detection is everything.

This is not an interpretation of $s$ but a prediction from it, and it can be
checked. Ordering the \fieldnum{155} hours of the budget window by their contribution to each term's variance gives Table~\ref{tab:shape_remedy}.

\begin{table}[t]
\caption{Concentration in time of each structural term over the hours of the budget window (clock-aligned basis of Table~\ref{tab:2}): shape statistic $s$; number (and fraction) of hours holding 90\,\% of the term's variance; the number of contiguous runs those hours form, with the number expected for the same count of hours drawn at random; share of the variance on the standard-deviation basis; and the remedy the shape indicates.}
\label{tab:shape_remedy}
\begin{tabular}{ll>{\raggedright\arraybackslash}p{0.17\textwidth}p{0.17\textwidth}p{0.12\textwidth}p{0.25\textwidth}}
\tophline
term & $s$ & hours holding 90\,\% of the variance & contiguous runs (random expectation) & share of variance (SD) & remedy \\
\middlehline
$I_0$ & \fieldnum{0.73} & \fieldnum{63} (\fieldnum{41}\,\%) & \fieldnum{24} (\fieldnum{38}) & \fieldnum{45}\,\% & longer calibration blocks (untested, Sect.~\ref{sec:6.2}) \\
$R$ & \textbf{\fieldnum{0.32}} & \textbf{\fieldnum{26} (\fieldnum{17}\,\%)} & \textbf{\fieldnum{12} (\fieldnum{22})} & \textbf{\fieldnum{52}\,\%} & detect and flag intervals (Sect.~\ref{sec:6.3}) \\
etalon $f$ & \fieldnum{0.98} & \fieldnum{118} (\fieldnum{76}\,\%) & \fieldnum{26} (\fieldnum{29}) & \fieldnum{3}\,\% & fit $f$ per scan, or model it \\
\bottomhline
\end{tabular}
\end{table}

The ordering by $s$ is the ordering by concentration, and the two-hour
interval from 23:00 UTC on 18 May to 01:00 UTC on 19 May alone holds \fieldnum{37}\,\% of the reflectivity term's variance.
The contiguous-run counts separate concentration from clumping: the
reflectivity hours form far fewer runs than the same number of hours drawn at
random, so they are episodes rather than scattered outliers and a flag can
address them. The etalon hours are indistinguishable from random, which is
the signature of a term no flag will help.

Read on the robust scale the ordering of the first two reverses -- $I_0$
carries \fieldnum{74}\,\% of the single-term quadrature sum and $R$ only \fieldnum{17}\,\% -- which is why the two scales
are reported side by side in Table~\ref{tab:2}. A term that is small for a typical
record and largest in total is exactly the one a campaign-mean uncertainty
hides.

\subsection{Calibration blocks: interval length and scans per block}
\label{sec:6.2}

The natural response to an interpolation error is to interpolate over shorter
intervals. How much that buys depends on the channel. Tripling the reflectivity interval from 2 to 6\,h was
measured to increase the interpolation error from \fieldnum{0.87}\,\% to \fieldnum{1.00}\,\% in hot PNs and from \fieldnum{2.05}\,\% to
\fieldnum{2.73}\,\% in hot ANs. In the 180\,$^{\circ}$C channel the error is set mostly by the noise on the knots, and a
shorter interval would buy little; in the 300\,$^{\circ}$C channel the interval carries curvature (Table~\ref{tab:splithalf_loo}),
and shorter intervals would help.

The other lever is the number of scans averaged into each calibration block, and the evidence on it is not
conclusive (Appendix~\ref{app:C}). Within the operational block ($N = 34$ scans, $T = 32$\,s), doubling the number of
scans averaged reduces the scatter of contiguous sub-blocks by \fieldnum{7}\,\% in the ANs channel, against the 29\,\% that
independent scans would give; but the range that can be measured this way lies entirely below the operational block
length, so it is mildly unfavourable rather than decisive, and the measurement that would settle it -- operating one
longer block -- has not been made. Even with white averaging the return would be modest: the zero-air term carries
\fieldnum{45}\,\% of the $\Sigma$ANs variance, so a fourfold block, which would raise the calibration duty from 1.20 to 3.89\,\%,
would lower the structural budget from \fieldnum{0.1236} to \fieldnum{0.1008}\,ppb (\fieldnum{18.5}\,\%), whereas flagging the few
defective reflectivity intervals would lower it to \fieldnum{0.0901}\,ppb. That ordering -- flagging before averaging --
holds whatever the block scaling turns out to be, because the reflectivity term is concentrated in a few episodes and
the zero-air term is not (Sect.~\ref{sec:6.1a}).

\subsection{Gating defective knots}
\label{sec:6.3}

Not all calibration blocks are usable, and the bad ones are few. Ordering the
cold-channel zero-air blocks by their leave-one-out influence, the worst
single block carries \fieldnum{28}\,\% of the total variance, the worst three carry \fieldnum{76}\,\%,
and the worst five carry \fieldnum{97}\,\%. A structural term computed over all blocks is
therefore a statement about a handful of them.

They are identifiable before any fit is attempted. A single observable -- the
block's light level relative to the campaign median -- with the threshold
$I/\tilde{I} < 0.5$ selects \fieldnum{13} of the \fieldnum{38} blocks whose influence exceeds 2\,\%,
but those \fieldnum{13} include nine of the worst ten and account for \fieldnum{99.9}\,\% of the
influence variance. The \fieldnum{25} blocks it misses are all in the \fieldnum{2}--\fieldnum{3}\,\% range. The
threshold excludes \fieldnum{8}\,\% of blocks.

Classification accuracy is the wrong figure of merit here and would have
rejected this threshold: \fieldnum{13} of \fieldnum{38} is a poor hit rate and a good filter. What
matters is the variance removed, and a single observable removes essentially
all of it. No multi-indicator scheme is needed.

The asymmetry with the reflectivity is worth recording, because it explains why
only one of the two references needed this. Reflectivity knots already pass
three quality gates -- a minimum fraction of valid pixels, an acceptance band on
the retrieved effective path length, and limits on the calibration residual --
so a defective helium or zero-air pair is rejected before it becomes a knot.
The zero-air reference had no equivalent: every block became a knot regardless
of how much light reached the detector. The gate described here closes that
gap. It is implemented with the threshold expressed as a fraction of the run's
own median block intensity rather than as an absolute count, so it transfers
between instruments, and it is disabled by default so that reprocessing an
existing product is an explicit choice rather than a silent change.

\subsection{Coverage: say when the reference is far away}
\label{sec:6.4}

Two things are missing, and conflating them produces the wrong prescription.
One is \emph{extrapolation}: records outside the range spanned by the reflectivity
knots. The other is a \emph{gap}: records between two knots more than two hours
apart. The hot channel has \fieldnum{3.76}\,\% of the first and \fieldnum{1.23}\,\% of the second over
the campaign; the cold channel \fieldnum{2.89}\,\% and \fieldnum{11.03}\,\%. Zero air has neither
(\fieldnum{0.08}\,\% and \fieldnum{0.07}\,\%), so the constraint is entirely on the reflectivity.

Two causes lie behind the deficit that survives clock alignment, and they cost different things to remove. The first
is bookkeeping and the cheapest fix in this paper: reflectivity generation stopped at 07-09, which produces the hot
channel's trailing extrapolation and is removed by re-running over files that already exist. The second is clustered
rejection of calibration pairs, \fieldnum{11.3}\,\% of cold pairs, which explains the cold gaps (\fieldnum{11.03}\,\% over the campaign and
\fieldnum{18.31}\,\% in the cold window); its remedy is to restore the light level or, failing that, redundancy. Only this
second cause is about how often the instrument calibrates, and even there the average rate is not the problem.
Independent rejection at the measured rate would leave \fieldnum{3.5}\,\% of cold records in gaps longer than two hours,
against \fieldnum{11.03}\,\% observed. The rejections arrive in runs: the worst cold gap, \fieldnum{11.8}\,h, is eleven consecutive
rejected pairs, which even at a 40\,\% independent rejection rate would be expected 0.03 times in the campaign. This
is the episodic structure Sect.~\ref{sec:6.1a} found in the reflectivity variance itself -- \fieldnum{26} of \fieldnum{155} hours carrying
90\,\% of it, in \fieldnum{12} contiguous runs where \fieldnum{22} would be expected at random -- seen from the calibration side.

It also locates the cost of quality control. The three gates the reflectivity already passes cost the hot channel
\fieldnum{0.24}\,\% of its pairs and the cold channel \fieldnum{11.3}\,\%; in the cold channel the gates are, correctly, the proximate
cause of the coverage deficit. This cuts against the argument of Sect.~\ref{sec:6.3}, which asks for a gate that zero
air does not yet have: gates remove variance and they remove coverage, and a pipeline that gates without also
reporting what the gating cost has moved the error somewhere it is harder to see. The context flag of
Sect.~\ref{sec:6.1} closes that loop. The same caution applies to the budget: the reflectivity contribution of
Sect.~\ref{sec:4} is measured by removing a knot and re-interpolating, so its magnitude scales with how far apart the
knots already are, and a sparse reference presents itself as a dominant term partly because it is sparse -- which is
why the remedy belongs to the schedule and to reprocessing rather than to the retrieval. Records in gaps are not
covered by the budget at all, since a perturbation experiment needs a knot to perturb; they are reported with the
context flag set and excluded from any statistic that assumes a characterised uncertainty.

In the heated channels zero air is measured every hour and helium every three, and each reflectivity knot pairs a
zero-air block with the most recent helium block that passed the quality gates (Appendix~\ref{app:A},
Table~\ref{tab:cal_schedule}), so the knots are hourly while their helium half is up to three hours old, and older when a
helium block is rejected. A rejected zero-air/helium pair removes one hourly knot and leaves a two-hour gap; a rejected
helium block removes no knot but makes the following knots use helium measured up to six hours earlier. Matching the
helium cadence to the zero-air cadence would bring the total calibration duty from 1.20\,\% to 1.79\,\% and make the
helium half of every knot contemporaneous. Whether it would also shorten the gaps depends on why pairs are rejected,
which in the cold channel is mainly the light level. The case for it is redundancy against runs of rejection and
contemporaneity, not precision: in the 180\,$^{\circ}$C channel the reflectivity error depends only weakly on interval
length (Sect.~\ref{sec:6.2}), so more frequent helium would not make a well-covered record much more accurate. At
1.2\,\% of the measurement time, calibration is the binding constraint on none of these remedies.

\subsection{Errors that neither the fit report nor the proposed fields detect}
\label{sec:6.5}

The convention of Sect.~\ref{sec:6.1} makes visible what the pipeline already knows. It cannot make visible what the
pipeline does not know, and two such errors appear in this record.

The first is a multiplicative gain error. Against a co-located NIER analyser the heated-channel NO$_2$ reads low by a
factor of about \fieldnum{0.6} while the ambient-temperature channel agrees (Sect.~\ref{sec:3.6}). A gain leaves the residual, the
reported uncertainty, the termination state and every distance and provenance field unchanged, and the zero-air test
of Sect.~\ref{sec:7.2} is blind to it because a gain applied to zero is zero. Only a comparison with an independent
measurement of the same quantity -- a reference analyser, a cylinder, or a channel that should agree -- reveals it,
and such a comparison does not by itself identify the cause.

The second is the raw light level. All six ambient-temperature excursions of Sect.~\ref{sec:7.3} coincide with spectra
that fell below 30\,\% of their usual count at 460\,nm; the residual finds them, and the reference context is abnormal
in five of them, but none of the proposed fields names the cause. The light level of each spectrum, and the helium and
zero-air intensities of each calibration block rather than only their ratio, are known when they are recorded and
could travel with the product as further context; we have not evaluated them for the heated channels.
